\section{An active microscopic realization}\label{sec:v13-active}
\subsection{Interventions and the unbounded flow}
Fix a nonempty finite system of $N\ge2$ identical hard spheres of
positive diameter on a flat torus, with positive configuration volume.
Let $\Gamma_N$ be the admissible
phase space and let $\mu$ be the normalized measure with density
proportional to
$\mathbf1_{\Gamma_N}(q,v)\exp(-b\sum_i|v_i|^2)$, $b>0$.
We use the classical almost-everywhere hard-sphere flow $\Phi_t$,
which is an invertible measure-preserving flow off a null exceptional
set; see \cite{GallagherSaintRaymondTexier} and the detailed
microscopic construction retained in the companion. No claim about
particle-number uniformity is needed. On $H=L^2(\mu)$,
\[
 U(t)f=f\circ\Phi_t
\]
is a strongly continuous unitary group, with skew-adjoint generator
$L$. Indeed invariance gives isometry and invertibility; for bounded
continuous phase functions, almost-everywhere continuity of the orbit
at zero and dominated convergence give strong continuity. Density and
isometry extend it to $H$. The generator domain includes the collision
identifications of the flow; it is not the unrestricted formal
transport derivative on arbitrary boundary data.

Let $A$ be finite and choose orthogonal matrices $O_a$ on velocity
space. The intervention
\[
 R_a(q_1,\ldots,q_N,v_1,\ldots,v_N)
    =(q_1,\ldots,q_N,O_av_1,\ldots,O_av_N)
\]
is an invertible $\mu$-preserving map. Put $C_af=f\circ R_a$. For a
fixed step length $\Delta>0$, the controlled physical step is
$F_a=\Phi_\Delta\circ R_a$, and its Koopman operator is
\begin{equation}\label{eq:v13-kick}
 T_a=C_aU(\Delta).
\end{equation}
The order matters. If both $O_a=I$ and $O_a=-I$ are allowed, the two
steps give different subsequent positions on a positive-measure set
of configurations with no collision in the next $\Delta$ units of
time and nonzero sufficiently small velocities. Such a set exists
inside any nonempty interior admissible configuration region by
choosing small velocities and then a positive interval away from zero.
The action therefore changes the dynamics, rather than merely choosing
which observable to read. We require no assertion that $C_aD(L)$ is
contained in $D(L)$.

For all finite control words, the flow and intervention compositions
are well defined on a common full-measure subset: the preimage of a
null exceptional flow set under each finite measure-preserving
composition is null, and there are countably many finite words.
Every preparation absolutely continuous with respect to $\mu$ is
supported there, including under a feedback choice of the word.

\subsection{A graph-core approximation and controlled words}
Choose an increasing family of finite-dimensional real subspaces
$V_n\subset D(L)$ whose union is a graph core, and let $P_n$ be the
orthogonal projections. Such a family can be constructed without
postulating smooth collision-boundary coordinates. Start with a
countable dense family $(f_j)$ in $H$ and a smooth compactly supported
approximate identity $\chi_\varepsilon$ on $\mathbb R$. Define
\[
 f_{j,\varepsilon}=\int_{\mathbb R}\chi_\varepsilon(s)U(s)f_j\,ds.
\]
A difference quotient followed by translation of $\chi_\varepsilon$
shows that $f_{j,\varepsilon}\in D(L^k)$ for every $k$ and
\begin{equation}\label{eq:v13-smoothing}
 L^kf_{j,\varepsilon}=(-1)^k\int
               \chi_\varepsilon^{(k)}(s)U(s)f_j\,ds,
 \qquad
 \|L^kf_{j,\varepsilon}\|_2
            \le\|\chi_\varepsilon^{(k)}\|_1\|f_j\|_2.
\end{equation}
For $f\in D(L)$ its smoothed versions converge to $f$ in graph norm:
commuting $L$ on its domain reduces this to strong continuity for $f$
and $Lf$. At each fixed smoothing scale, approximation of $f$ by the
dense family controls both the function and its first derivative by
\eqref{eq:v13-smoothing}. A countable sequence of scales and finite
linear spans consequently gives a graph core. Enumerate those vectors
and let $V_n$ be their successive spans. This also shows that their
collision compatibility comes from the actual flow domain, not from
an imposed incompatible trace.

Set $L_n=P_nLP_n$ on $V_n$, $U_n(t)=\exp(tL_n)$ on $V_n$, and, as an
operator on $H$ extended by zero on $V_n^\perp$, define
\begin{equation}\label{eq:v13-approxword}
 T_{a,n}=P_nC_aP_nU_n(\Delta)P_n.
\end{equation}
These operators are contractions. Their finite-dimensional matrices
need not be Markov matrices: they approximate \emph{observables}; the
observation experiment constructed below is nevertheless a genuine
probability experiment.

\begin{lemma}[Graph-core and controlled-word convergence]
\label{lem:v13-words}
For every $f\in H$ and $T_0<\infty$,
\begin{equation}\label{eq:v13-orbit}
 \sup_{0\le t\le T_0}\|U_n(t)P_nf-U(t)f\|_2\longrightarrow0.
\end{equation}
For every fixed word length $k$,
\begin{equation}\label{eq:v13-words}
 \max_{(a_1,\ldots,a_k)\in A^k}
 \|T_{a_1,n}\cdots T_{a_k,n}f-T_{a_1}\cdots T_{a_k}f\|_2
                         \longrightarrow0.
\end{equation}
Neither statement needs domain invariance under the interventions.
\end{lemma}
\begin{proof}
The matrix $L_n$ is skew-adjoint on $V_n$ because $L$ is
skew-adjoint and $V_n\subset D(L)$. Thus its group is unitary and
$\|(z-L_n)^{-1}\|\le1/\Re z$ when $\Re z>0$.
If $u\in V_m$ and $n\ge m$, then
\[
 (z-L_n)^{-1}P_n(z-L)u=u.
\]
The set $(z-L)\bigcup_mV_m$ is dense in $H$, since the union is a
graph core and $z-L$ maps $D(L)$ onto $H$. The uniform resolvent
bounds and $P_n\to I$ strongly therefore give
\[
 (z-L_n)^{-1}P_nf\longrightarrow(z-L)^{-1}f\quad(f\in H).
\]
This is the graph-core consistency and stable-resolvent setting of
Trotter's theorem \cite[Theorem~5.2]{Trotter1958}. We give the
contraction argument to specify the hypotheses used here.
For any contraction generator $A$ and $u\in D(A)$, the Euler
resolvent formula can be written as a Gamma integral:
\[
 (I-tA/m)^{-m}u=\int_0^\infty e^{sA}u\,d\gamma_{m,t}(s),
 \qquad \E_{\gamma_{m,t}}s=t,\quad
               \Var_{\gamma_{m,t}}s=t^2/m.
\]
This follows by multiplying the Laplace resolvent integrals and
convolving their exponential densities. Since
$\|e^{sA}u-e^{tA}u\|\le|s-t|\|Au\|$, its error from $e^{tA}u$
is at most $t\|Au\|/\sqrt m$.
For $u\in V_l$ one has $\|L_nu\|\le\|Lu\|$ for all $n\ge l$.
At each fixed $t>0$ and Euler order $m$, the resolvent convergence
already proved gives convergence of its $m$th power. First letting
$n\to\infty$ and then $m\to\infty$ proves orbit convergence on the
core. The orbit Lipschitz bound with constant $\|Lu\|$, uniform in
$n$, upgrades it to uniformity on $[0,T_0]$ by a finite time mesh.
Density and contraction extend it to all $f\in H$. No derivative of
$U(t)f$ for $f\notin D(L)$ was taken.

Now $P_nC_aP_n\to C_a$ strongly by boundedness of $C_a$ and strong
convergence of $P_n$. Equation~\eqref{eq:v13-orbit} gives
$T_{a,n}\to T_a$ strongly. For a fixed finite word, subtract the two
products one factor at a time. Each difference acts on a fixed vector
formed from the limiting factors, while every factor to its left is
a contraction. Each term tends to zero. There are finitely many words
of length $k$, which proves the maximum in \eqref{eq:v13-words}.
At no step was an expression such as $LC_a f$ required.
\end{proof}

\subsection{Actual marked acquisition under feedback}
Fix real observables $f_k\in H$ and noises $\sigma_k>0$. At step $k$,
a controller chooses $a_k$ from its previously available reports and
retained state, the action is applied to the physical state, and the
new report is
\begin{equation}\label{eq:v13-active-report}
 Y_k=(T_{a_1}\cdots T_{a_k}f_k)(x_0)+\sigma_k\xi_k,
 \qquad \xi_k\text{ independent standard Gaussians}.
\end{equation}
The initial law is $p_\theta\,d\mu$, with
$\sup_\theta\|p_\theta\|_2\le D<\infty$.
An unobserved initial physical mark $W$ is drawn jointly with $x_0$
through a fixed channel $\gamma(dw\mid x_0)$, independently of the
acquisition noises. Losses may use that actual mark and the observed
control--report transcript. This specification does not allow an
unproved replacement of a controlled terminal physical quantity by
an independent variable of the same marginal law.

The approximation replaces the product in
\eqref{eq:v13-active-report} by $T_{a_1,n}\cdots T_{a_k,n}f_k$.
It has a finite-dimensional physical realization. In a real orthonormal
basis $(\varphi_1,\ldots,\varphi_{d_n})$ of $V_n$, write the matrix of
$T_{a,n}$ as $M_{a,n}$ and the coefficient vector of $P_nf_k$ as
$c_{k,n}$. The initial row vector is
$v_0=(\varphi_1(x_0),\ldots,\varphi_{d_n}(x_0))$, and
\[
 v_k=v_{k-1}M_{a_k,n},\qquad Y_k^{(n)}=v_kc_{k,n}+\sigma_k\xi_k.
\]
Initialize $(v_0,W)$ with the actual push-forward joint law of $(x_0,W)$.
This is a valid controlled probability model with a real vector state;
no positivity of $M_{a,n}$ is needed. Its real dimension $d_n$ is not
counted as $d_n$ finite retained labels. The consumer register remains
a separate constrained object. Measurable representatives of the
finitely many functions may be chosen once; their null ambiguities
have zero probability under the specified preparations.

\begin{theorem}[Active marked deficiency bound]\label{thm:v13-active}
For each fixed horizon $T$, define
\begin{equation}\label{eq:v13-active-error}
 \epsilon_{n,T}=1\wedge\frac{D}{\sqrt{2\pi}}
       \sum_{k=1}^T\frac1{\sigma_k}
       \sum_{a\in A^k}
       \|T_{a_1}\cdots T_{a_k}f_k
                     -T_{a_1,n}\cdots T_{a_k,n}f_k\|_2.
\end{equation}
The exact and approximate acquisitions have two-sided stateless marked
causal deficiency at most $\epsilon_{n,T}$, using the identity report
channel. The bound includes every independent visible seed in the
joint comparison, holds uniformly over every finite-register consumer
and feedback design, and tends to zero for each fixed $T$.
Consequently all paired $[0,1]$ task risks and optimized lower support
values have error at most $\epsilon_{n,T}$ at unchanged widths; paired
centered values have error at most $2\epsilon_{n,T}$ when their
unrestricted benchmarks are also paired.
\end{theorem}
\begin{proof}
Couple the two experiments with the same $x_0,W$ and every independent
visible seed. Couple their private controller randomness so that the
same available history produces the same action. At each step before
the first report mismatch, use a maximal coupling of the two Gaussian
reports. Gaussian densities satisfy
\[
 \|N(u,\sigma^2)-N(v,\sigma^2)\|_{\rm TV}
     =2\Phi(|u-v|/(2\sigma))-1
     \le\frac{|u-v|}{\sigma\sqrt{2\pi}}.
\]
The equality follows by separating the densities at their midpoint;
the inequality follows by bounding the standard normal density by
$1/\sqrt{2\pi}$. Maximal couplings can be chosen measurably from
these densities by their common part and residual parts.

For a fixed $x_0$, the probability of a first mismatch at time $k$ is
bounded by the sum, over the possible matched action words $a\in A^k$,
of their matched-history probabilities times the displayed Gaussian
bound for that word. Discarding these probabilities gives an upper
bound by the sum over all words. Integrate under the \emph{original}
preparation $p_\theta\,d\mu$. Cauchy--Schwarz bounds the integral of
each absolute mean difference by $D$ times its $L^2$ norm. Summing
first-mismatch probabilities yields \eqref{eq:v13-active-error}.
In particular no bound on a feedback-conditioned posterior density
has been assumed. When no mismatch occurs the entire transcript and
actual mark, including the visible seed, coincide. The coupling
inequality gives their joint total variation bound.

The same argument is symmetric in the two report models; the
stateless identity channel is executable in both directions, so it
bounds the intrinsic infima. Lemma~\ref{lem:v13-words} makes every term
in the fixed finite sum tend to zero. The comparison precedes any
consumer optimization, hence its bound does not depend on the
consumer's number of labels. The remaining conclusions are the
channel and support inequalities in
Theorem~\ref{thm:v13-main} and \eqref{eq:v13-risktransport}, which
remain valid on these standard Borel carriers.
\end{proof}

\subsection{Active regenerative average decisions}
Take a fixed observable $f$, fixed noise $\sigma>0$, and a fixed finite
measurable report partition applied to each Gaussian acquisition.
At each cycle boundary the physical system is re-prepared from
$p_\theta\,d\mu$ with its actual mark channel $\gamma$, and the
within-policy register is reset to $m_0$. Between boundaries the
orthogonal interventions and hard-sphere flow continue as above. The
independent reset probability is $\eta>0$. In the approximate system
reset the real vector and mark by the same joint push-forward
preparation. A stationary consumer uses the partitioned reports,
actions, and its finite register. Its $[0,1]$ one-stage losses depend
on the actual cycle mark, control, and current finite report; reset
stage costs are fixed and identically paired. A finite task index is
available to the controller but not to the common encoder.

Assume $j\mapsto p_j$ is continuous in $L^1(\mu)$ on a compact row
set, the $L^2$ bounds are uniform, and the bounded row losses vary
uniformly continuously. The flow, intervention maps, observable,
noise law, report partition, and mark channel are fixed across rows.
Then the exact and every approximate model satisfy
\eqref{eq:v13-continuity}: only the reset law and losses vary with the
row, and the dependence on finite policy rows is a finite polynomial
expansion in kernels of norm one. The underlying physical transitions
are allowed to be deterministic between resets.

\begin{corollary}[Active average-risk realization]\label{cor:v13-active-average}
Both models have the compact average-risk bodies and attained duality
of Theorem~\ref{thm:v13-average}. At every finite consumer width their
paired average risks, lower supports, and equally centered minimax
values differ by at most
\begin{equation}\label{eq:v13-active-average}
 \mathfrak e_n=\eta\sum_{k=0}^\infty
                    (1-\eta)^k\epsilon_{n,k+1}\longrightarrow0.
\end{equation}
The convergence is uniform in the consumer width and in any independent
priced mixture of stationary designs. It does not require a controlled
word approximation rate uniform in word length.
\end{corollary}
\begin{proof}
A deterministic report partition contracts total variation.
The age-$k$ stage is evaluated after at most $k+1$ controlled
acquisitions, so Theorem~\ref{thm:v13-active} supplies the cycle
comparison with $\varepsilon_{n,k}=\epsilon_{n,k+1}$. Apply
Theorem~\ref{thm:v13-cycle}. Its geometric tail proof gives the limit
although the finite-word bound in \eqref{eq:v13-active-error} can grow
with $k$. The reset and row-continuity assertions just verified give
Theorem~\ref{thm:v13-average}. Conditioning on each persistent mode
and charging its index proves the mixture statement.
\end{proof}

\begin{remark}[Fixed physical scope]
The interventions above are active changes to microscopic dynamics,
and the average theorem uses actual repeated physical preparation.
Neither observation control alone nor an assumed posterior $L^2$
bound is substituted for them. The assertions are for fixed $N$,
sphere diameter, $A$, $\Delta$, and positive noise, and for the declared
initial-mark losses. No estimate is asserted uniform in particle
number, a growing action alphabet, vanishing noise, or a kinetic
scaling. The historical nonlinear kinetic corrector and the Sinai
left-strip spectral problem remain different analytic statements.
They are preserved in the development with their original hypotheses,
not inferred from the present operator approximation.
\end{remark}
